\begin{definition}
	Пусть $E$ "--- линейное нормированное пространство, $L \subset E$ "--- линейное нормированное подпространство, $h \in E$. \textit{Элементом наилучшего приближения} для $h$ называется $x \in L$ такой, что 
    $\|h - x\| = \rho(h, L) = \inf \limits_{y \in L} \|h - y\|$.
\end{definition}

\begin{note}
	Можно построить примеры, в которых элемент наилучшего приближения не существует или не единственен.
\end{note}

\begin{example}
Рассмотрим $\R^2_{\infty}$, $\|x\| = \max(|x_1|, |x_2|)$.
Пусть $L = \{x : x_2 = 0\}$, $x = (1, 1)$. Так как открытый шар 
в $\R^2_{\infty}$ выглядит как квадрат, то открытый шар
с центром в $x$ коснётся $L$ сразу во многих точках
(а именно, в $\{ (x_1, 0) : 0 \leqslant x_1 \leqslant 2 \})$.
\end{example}

\begin{proposition}\label{bestelt}
	Пусть $H$ "--- гильбертово пространство, $M \subset H$ "--- подпространство в $H$. Тогда для любого $h \in H$ существует единственный элемент наилучшего приближения $x \in M$.
\end{proposition}

\begin{proof}
	Зафиксируем $h \in H$. Сначала докажем, что элемент наилучшего приближения для $h$ существует. Положим $d := \rho(h, M)$. Если $d = 0$, то, в силу замкнутости множества $M$, выполнено $h \in M$, и в качестве элемента наилучшего приближения для $h$ подходит сам $h$. Иначе --- выберем $\{x_n\} \subset M$ такую, что $\|h - x_n\| \to d$. По равенству параллелограмма, для любых $n, m \in \N$ выполнено следующее:
	\[\|x_n - x_m\|^2 = 2\|h - x_n\|^2 + 2\|h - x_m\|^2 - 4\left\|h - \frac{x_n + x_m}{2}\right\|^2\]

	Поскольку $\|h - x_n\|^2 \to d^2$ при $n \to \infty$, $\|h - x_m\|^2 \to d^2$ при $m \to \infty$ и выполнено неравенство $\|h - \frac{x_n + x_m}{2}\|^2 \ge d^2$, последовательность $\{x_n\}$ фундаментальна. Поскольку пространство $H$ полно, а $M$ замкнуто, существует $x \in M$ такой, что $x_n \xrightarrow[H]{} x$, причем, в силу непрерывности нормы, $\|h - x\| = d$.

	Покажем теперь, что элемент наилучшего приближения для $h$ единственен. Пусть для некоторого $y \in M$ тоже выполнено равенство $\|h - y\| = d$, тогда, по равенству параллелограмма, выполнено следующее $\left[h - x + h - y = 2h - x -y = 2(h - \frac{x+y}{2}), \; h - x - h + y = -(x - y)\right]$:
	\[4d^2 = 2\|h - x\|^2 + 2\|h - y\|^2 = \|x - y\|^2 + 4\left\|h - \frac{x+y}{2}\right\|^2 \geqslant \|x - y\|^2 + 4d^2\]

	Значит, $\|x - y\| = 0$, то есть $x = y$, и элемент наилучшего приближения единственен.
\end{proof}


\begin{definition}
	Пусть $E$ "--- евклидово пространство, $S \subset E$. \textit{Аннулятором} множества $S$ называется следующее множество:
	\[S^\perp := \{y \in E: \forall x \in S: (x, y) = 0\}\]
\end{definition}

\begin{note}
	Легко проверить, что $S^\perp$ является подпространством в $E$. Кроме того, выполнены равенства $S^\perp = [S]^\perp = \overline{[S]}^\perp$.
\end{note}

\begin{theorem}[Рисса, о проекции]\label{thm4.3}
	Пусть $H$ "--- гильбертово пространство, {$M \subset H$} -- подпространство в $H$. Тогда $H = M \oplus M^\perp$.
\end{theorem}

\begin{proof}
	Покажем сначала, что $H = M + M^\perp$. Зафиксируем $h \in H$, и по утверждению \ref{bestelt} выберем его элемент наилучшего приближения $x \in M$. Положим $y := h - x$ и $d := \|y\|$, тогда для произвольного $m \in M$ и произвольного $\alpha \in \R \bs \{0\}$ выполнено следующее:
	\[d^2 = \|h - x\|^2 \leqslant \|h - (x + \alpha m)\|^2 = d^2 - 2\alpha(h - x, m) + \alpha^2\|m\|^2\]

	Следовательно, $(y, m) \leqslant \frac\alpha2\|m\|^2$ для любого $\alpha \in \R \bs \{0\}$, откуда $(y, m) = 0$. В силу произвольности выбора $m \in M$, получаем, что $y \in M^\perp$, причем $h = x + y$. Значит, выполнено равенство $H = M + M^\perp$. Проверим теперь, что сумма $M + M^\perp$ --- прямая. Действительно, если $z \in M \cap M^\perp$, то $(z, z) = 0$, откуда $z = 0$, что и означает требуемое.
\end{proof}

\begin{note}
	Теорема выше отражает одно из важных свойств гильбертовых пространств, которое для банаховых пространств верно не всегда. Далее мы обсудим еще одно такое свойство --- существование базиса в бесконечномерном сепарабельном пространстве.
\end{note}